\documentclass[12pt,letterpaper,onecolumn]{article}
\usepackage[T1]{fontenc}
\usepackage{amsmath,amssymb,amsthm,bm,mathtools}
\usepackage{newtxtext,newtxmath}
\usepackage{booktabs,array,microtype}
\usepackage{cite}

\usepackage{xcolor,url}
\usepackage[unicode,bookmarksnumbered,breaklinks,colorlinks=false,
 citebordercolor={0 1 0},linkbordercolor={1 0 0},urlbordercolor={0 1 1},
 pdfborder={0 0 0.65}]{hyperref}
\newcommand{\doi}[1]{\href{https://doi.org/#1}{\nolinkurl{#1}}}
\newcommand{\arxiv}[1]{\href{https://arxiv.org/abs/#1}{arXiv:\nolinkurl{#1}}}
\allowdisplaybreaks[2]
\interdisplaylinepenalty=2500
\setlength{\emergencystretch}{2em}

\newtheorem{theorem}{Theorem}
\newtheorem{lemma}[theorem]{Lemma}
\newtheorem{proposition}[theorem]{Proposition}
\newtheorem{corollary}[theorem]{Corollary}
\theoremstyle{definition}
\newtheorem{definition}[theorem]{Definition}
\theoremstyle{remark}
\newtheorem{remark}[theorem]{Remark}
\DeclareMathOperator{\Tr}{Tr}
\DeclareMathOperator{\supp}{supp}
\DeclareMathOperator{\rank}{rank}
\DeclareMathOperator{\diag}{diag}
\DeclareMathOperator{\conv}{conv}
\DeclareMathOperator{\spec}{spec}
\newcommand{\C}{\mathbb C}
\newcommand{\R}{\mathbb R}
\newcommand{\Z}{\mathcal Z}
\newcommand{\G}{\mathbf G}
\newcommand{\T}{\mathbf T}
\newcommand{\Id}{\mathbf I}
\newcommand{\Pmat}{\mathbf P}
\newcommand{\Bcl}{B_{\mathrm{ABL},q}}
\newcommand{\CA}{C_{\mathrm A,q}}
\newcommand{\hbin}{h}
\newcommand{\ket}[1]{\lvert #1\rangle}
\newcommand{\bra}[1]{\langle #1\rvert}
\newcommand{\ip}[2]{\langle #1,#2\rangle}
\newcommand{\HS}{\mathrm{HS}}
\newcommand{\Rmc}{B_{\mathrm{MC},q}}
\newcommand{\Bmix}{B_{\mathrm{mix},q}}
\newcommand{\Bcomm}{B_{\mathrm{comm},q}}
\newcommand{\rp}{r_{\pi}}
\newcommand{\Dperp}{D_{\perp,q}}
\newcommand{\Dpar}{D_{\parallel,q}}
\hyphenation{non-bi-nary Hole-vo Aal-to-nen}




\usepackage[margin=1in]{geometry}
\hypersetup{colorlinks=true,linkcolor=black,citecolor=black,urlcolor=black,pdfauthor={Yang Shifu},pdfsubject={Mathematical proof note},pdfcreator={LaTeX}}
\pagestyle{plain}
\renewcommand{\refname}{Sources}

\begin{document}
\noindent{\large\bfseries Mathematical proof note}\par\medskip
\noindent{\raggedright\fontsize{24}{28}\selectfont\bfseries Strict Nonbinary Rate Bounds Beyond Optimized Classical and Mixed-Channel Envelopes\par}\medskip
\noindent{\fontsize{16}{20}\selectfont Yang Shifu}\par\smallskip
\noindent\begin{minipage}{\linewidth}\small National Mobile Communications Research Laboratory and the Frontiers Science Center for Mobile Information Communication and Security, Southeast University, Nanjing 211189, China\end{minipage}\par\medskip
\noindent Version 1.0 --- October 8, 2026\par\bigskip
\hypersetup{pdftitle={Strict Nonbinary Rate Bounds Beyond Optimized Classical and Mixed-Channel Envelopes}}
\section*{Result overview}
For every fixed integer alphabet size \(q\ge3\), we strictly improve the joint flagged envelope of the full-domain Aaltonen/product/shortening family specified here and the unmasked \(q\)-ary mixed-channel family near the Plotkin endpoint. Cyclic Gram channels represent the classical envelope exactly. Conditional support expansion improves boundary generators, while independent coherence variations improve full-rank mixed-channel generators outside a finite stationary set. A positive information slope for the remaining generators forces every sufficiently low-rate optimizer to contain an improvable component. The improved channels have output dimension at most \(2q^2+q\), and the resulting rate bound applies to linear and nonlinear codes.
\section{Statement and Auxiliary-Channel Criterion}
Let \(q\ge3\) be a fixed integer, \(k=q-1\), and \(\theta=k/q\). For the maximum cardinality \(A_q(n,d)\) of an arbitrary \(q\)-ary code of minimum Hamming distance at least \(d\), put
\[
R_q(\delta)=\limsup_{n\to\infty}n^{-1}\log_q A_q(n,\lceil\delta n\rceil).
\]
The benchmark jointly optimizes all parameters of the named classical components, all unmasked \(q\)-MC inputs, and all finite orthogonal flags.

For a finite cyclically output-symmetric channel \(x\mapsto\sigma_x\), write
\begin{align}
P&=q^{-1}\sum_x\sigma_x,\notag\\
\chi&=S(P)-q^{-1}\sum_xS(\sigma_x),\notag\\
p&=1-q^{-2}\sum_x
\Tr(P^{-1/2}\sigma_xP^{-1/2}\sigma_x).\label{eq:cq}
\end{align}
Here \(S(A)=-\Tr A\ln A\), and inverses act on mean support. Cyclic symmetry means \(U\sigma_xU^*=\sigma_{x+1}\), with indices modulo \(q\). The \(q\)-ary pretty-good-measurement criterion~\cite[Prop.~27]{AG2026} gives
\begin{equation}\label{eq:criterion}
p<\delta\quad\Longrightarrow\quad R_q(\delta)\le\chi/\ln q.
\end{equation}
The channel is fixed before the ordinary blocklength limit. For input-independent orthogonal flags,
\(\sigma_x=\bigoplus_i a_i\sigma_{i,x}\), both \(p\) and \(\chi\) are affine in the weights \(a_i\).

Sections~\ref{sec:classical} and~\ref{sec:MC} define two compact generator families \(\mathfrak A_q,\mathfrak M_q\). Define their full flagged envelope by
\begin{equation}\label{eq:benchmark}
\begin{gathered}
B_q(\delta)=\min_{\substack{(p,r)\in\conv\mathcal K_q\\p\le\delta}}r,
\\
\mathcal K_q=\{(p_\sigma,\chi_\sigma/\ln q):
\sigma\in\mathfrak A_q\cup\mathfrak M_q\}.
\end{gathered}
\end{equation}

\begin{theorem}\label{thm:main}
For every fixed integer \(q\ge3\), there is \(\eta_q>0\) such that, for every
\(\theta-\eta_q\le\delta<\theta\), a cyclic channel \(\rho\) of output dimension at most \(2q^2+q\) satisfies
\begin{equation}\label{eq:strict}
p_\rho<\delta,\qquad
R_q(\delta)\le\chi_\rho/\ln q<B_q(\delta).
\end{equation}
At each interior point of this interval, one fixed channel retains a uniform positive gap on a compact neighborhood.
\end{theorem}

The proof gives \(\eta_q=\theta(\beta_q/(4q))^2\), with a strictly positive \(\beta_q\) defined from a finite stationary set in~\eqref{eq:beta}. No uniformity in growing \(q\) is asserted.

\section{The Complete Named Classical Envelope}\label{sec:classical}
We first specify the comparison domain. Put
\[
H_q(x)=\frac{-x\ln x-(1-x)\ln(1-x)+x\ln k}{\ln q}.
\]
Aaltonen's scalar source~\cite[Eq.~(9) and concluding theorem]{Aaltonen1990} has triples \((\omega,\xi,\eta)\) satisfying
\begin{align}
0\le\omega\le1,\quad
0\le\eta&\le\frac{q-2}{q-1}\omega,\notag\\
0\le\xi-\eta&\le\min(\omega-\eta,1-\omega),\label{eq:domain}\\
\beta_*&\le\omega-\frac{q-1}{q-2}\eta.\notag
\end{align}
For \(s=(\omega-\eta)/(1-\eta)\) and \(t=(\xi-\eta)/(1-\eta)\), the threshold and objective are
\begin{align}
\beta_*&=(1-\eta)\frac{s(1-s)-t(1-t)}{1+2\sqrt{t(1-t)}},\notag\\
e_A&=2\beta_*+(\omega-\beta_*)k_{q-1}
                 \left(\frac{\eta}{\omega-\beta_*}\right),\label{eq:source}\\
f_q(\xi,\eta)&=\frac{H(\eta,\xi-\eta,1-\xi)+\eta\ln(q-2)}{\ln q},\notag\\
k_m(x)&=\frac{m-1-(m-2)x-2\sqrt{(m-1)x(1-x)}}m.\notag
\end{align}
Here \(H\) is Shannon entropy in natural units; collapsed expressions use continuous values. The whole-space objective is \(1-H_q(\omega)+f_q(\xi,\eta)\), with \(e_A\le\delta\).
The original free parameter \(\beta\ge\beta_*\) can be eliminated without loss: the derivative of
\(2\beta+(\omega-\beta)k_{q-1}(\eta/(\omega-\beta))\) is
\[
\frac{q+\sqrt{(q-2)\eta/(\omega-\beta-\eta)}}{q-1}>0.
\]

\subsection{A compact Gram parameterization}
Let \(e_0\) be the first basis vector, \(e_+=k^{-1/2}\sum_{j>0}e_j\), and \(Q_\perp=I-e_0e_0^*-e_+e_+^*\). Use the compact domain
\[
\mathcal Z=\{(s,u,v):0\le s\le1,\ 0\le u\le\sqrt{s(1-s)},\
0\le v\le s+u\}.
\]
For \(A=s(1-s)-u^2\) and \(D=1+2u+(q-2)v^2\), set
\begin{equation}\label{eq:T}
T=D^{-1/2}\left[
vQ_\perp+
\begin{pmatrix}1-s+u&\sqrt A\\\sqrt A&s+u\end{pmatrix}_{e_0,e_+}
\right],\qquad G=T^2.
\end{equation}
The two-dimensional block has determinant \(u(1+2u)\) before normalization and nonnegative diagonal, so \(T\succeq0\). Expansion gives \(\Tr T^2=1\).

\begin{lemma}\label{lem:source}
Every source triple in~\eqref{eq:domain} is represented by~\eqref{eq:T}, and conversely. With
\begin{align}
\eta&=(q-2)v^2/D,& \alpha&=(1+2u)/D,\notag\\
t&=(1-\sqrt{1-4u^2})/2,&
\omega&=\eta+\alpha s,\quad \xi=\eta+\alpha t,\label{eq:map}
\end{align}
the source pair is
\begin{align}
F&=\ln q-H(\diag G)+S(G),\label{eq:F}\\
p&=1-\sum_jT_{jj}^2
=\frac{2A+(q-2)(s+u-v)^2/(q-1)}D=e_A.\label{eq:p}
\end{align}
Its rate coordinate is \(F/\ln q=1-H_q(\omega)+f_q(\xi,\eta)\).
\end{lemma}
\begin{proof}
Squaring the block gives
\begin{gather*}
\diag G=(1-\omega,\omega/k,\ldots,\omega/k),\\
\spec G=\{\alpha(1-t),\alpha t,(\eta/(q-2))^{[q-2]}\}.
\end{gather*}
These identities prove~\eqref{eq:F}. Moreover
\[
T_{00}=(1-s+u)/\sqrt D,\qquad
T_{jj}=\frac{s+u+(q-2)v}{(q-1)\sqrt D}\quad(j>0),
\]
giving~\eqref{eq:p}. For the inverse use
\begin{align*}
s&=\frac{\omega-\eta}{1-\eta},&
t&=\frac{\xi-\eta}{1-\eta},\\
u&=\sqrt{t(1-t)},&
v&=\sqrt{\frac{\eta(1+2u)}{(q-2)(1-\eta)}}.
\end{align*}
The degree constraint is exactly \(u^2\le s(1-s)\).
Also
\[
\beta_*=A/D,\quad
\omega-\frac{q-1}{q-2}\eta=\frac{(1+2u)s-v^2}{D},
\]
so the last source constraint is \(v\le s+u\).
It implies the remaining constraint because
\(v^2\le(s+u)^2\le s(1+2u)\).
Finally \(\omega-\beta_*-\eta=(s+u)^2/D\); insertion into~\eqref{eq:source} proves its threshold identity. Here \(\eta<1\) throughout the source domain, and all zero-prior endpoints are retained.
\end{proof}

Writing \(v_j=Te_j\), define the classical cyclic generator
\begin{equation}\label{eq:lift}
\sigma_x=\bigoplus_{c=0}^{q-1}v_{x+c}v_{x+c}^*.
\end{equation}
This is a density matrix in dimension \(q^2\), and shifting the blocks implements cyclic symmetry. Its mean is \(\bigoplus_cG/q\); the output eigenvalues are \(\diag G\), so its information is \(F\). Its PGM success is
\(\sum_j(v_j^*G^{-1/2}v_j)^2=\sum_jT_{jj}^2\), since \(TG^{-1/2}T=T\) on support. Thus~\eqref{eq:lift} realizes the exact source pair, even for singular \(G\). Let \(\mathfrak A_q\) contain all these generators.

\subsection{Products, shortening, and endpoint components}
Let \(W_q(\omega,d)\) be the infimum of \(f_q\) at fixed \(\omega\) and \(e_A\le d\), with the included zero-rate constant-weight branch at
\(d\ge D_0(\omega)=2\omega-q\omega^2/(q-1)\).
Define \(B_{\rm cl}\) as the closure of the source bound under the full two-block construction~\cite[Thms.~6--7]{BHL2007}, ordinary shortening, and the Laihonen--Litsyn operation~\cite[Thm.~1]{LL1998}
\begin{equation}\label{eq:LL}
\tau[1-H_q(w)]+(1-\tau)B(d'),\quad
\delta=2\tau w+(1-\tau)d'.
\end{equation}
Here \(0\le w\le\theta\), \(0<\tau<1\), and \(0\le d'\le1\). Include the Hamming, Elias, Singleton, Plotkin and first-LP components, finite recursive uses with their extraction terms, and their closed limiting values.

\begin{lemma}\label{lem:closure}
This named classical closure is exactly the lower error-budget boundary of the convex hull of the pairs of \(\mathfrak A_q\).
\end{lemma}
\begin{proof}
The rank-one point \((s,u,v)=(\omega,0,0)\) has pair
\((D_0(\omega),1-H_q(\omega))\); its actual error is at most the Hamming charge \(2\omega\). The endpoints \((0,1)\) and \((\theta,0)\) occur in the family. The first-LP pair is also included: for \(f=q^{-1/2}\sum_je_j\), use
\[
G_x=\frac{x}{k}(I-ff^*)+(1-x)ff^*,\qquad 0\le x\le\theta.
\]
It corresponds to \(\omega=\theta,\eta=(q-2)x/k,\xi=x\).
Its information is \((\ln q)H_q(x)\) and its error is
\[
\gamma_q(x)=1-\frac{(\sqrt{1-x}+\sqrt{kx})^2}{q}.
\]
The source constraints follow from \(x\le\theta\) and
\[
\beta_*=\frac{\theta}{q}
\left(\sqrt{1-x}-\sqrt{x/k}\right)^2\le\theta-x,
\]
using \((\sqrt a-\sqrt b)^2\le a-b\) for \(a\ge b\).
Since \(\gamma_q\) is an involution on \([0,\theta]\), this gives \(H_q(\gamma_q(\delta))\).

For the two-block construction, the constant-weight bound is
\(\nu W_q(\omega_1,d_1)+(1-\nu)W_q(\omega_2,d_2)\).
The extraction terms are
\(H_q(\omega)-\nu H_q(\omega_1)-(1-\nu)H_q(\omega_2)\) and \(1-H_q(\omega)\). They cancel to give
\begin{align}\label{eq:product}
&\nu[1-H_q(\omega_1)+W_q(\omega_1,d_1)]\notag\\
&\qquad +(1-\nu)[1-H_q(\omega_2)+W_q(\omega_2,d_2)].
\end{align}
The constraints are \(\omega=\nu\omega_1+(1-\nu)\omega_2\),
\(\delta=\nu d_1+(1-\nu)d_2\), and
\(0\le d_i\le\min(2\omega_i,1)\).
Each component is a source pair with actual error at most its charged distance, so a flag realizes a pair no worse than~\eqref{eq:product}.
The same argument treats~\eqref{eq:LL} and ordinary shortening. The endpoint chord \(1-\delta/\theta\) dominates Singleton. Rank-one points handle Elias and Hamming. The LL first-LP straight line is the shortening chord through
\(\delta=(q-2)^2/q^2\), where \(\gamma_q(\delta)=1/q\).
Induction and compactness handle recursion and limits.

For the reverse inclusion, a lower-hull minimizer at an interior error uses at most two source points. Indeed zero information occurs only at \((\theta,0)\); hence the minimum is positive and has actual error \(\delta\), and an exposed supporting face has dimension at most one. Given the two compositions and errors, set \(\varphi=\nu\omega_1\) in the original product domain and take \(d_i=e_i\). Its constraints hold because \(e_i\le\theta\) and \(e_i\le2\omega_i\). The latter follows by expanding~\eqref{eq:source}:
\[
e_A=2\omega-\frac{qy+(q-3)\eta+
2\sqrt{(q-2)\eta(y-\eta)}}{q-1},\quad y=\omega-\beta_*.
\]
Each selected source is feasible for \(W_q(\omega_i,e_i)\), so~\eqref{eq:product} attains a value no larger than the convex boundary. Endpoints follow by continuity.
\end{proof}

The closed zero-rate and shortening conventions used here incur only a polynomial endpoint correction. At \(d=nD_0(\omega)\), centered simplex images of constant-weight words have pairwise nonpositive inner products in dimension \(n(q-1)\), hence at most \(2n(q-1)\) words: their unit Gram matrix \(K\) has \(\Tr K^2\le2\Tr K\), so \(\rank K\ge |C|/2\). The zero-norm case has one word. Small positive source degrees approach the same zero-rate value from below the threshold. For LL ball selection at equality \(d=2r>0\), a puncturing fiber has disjoint weight-\(r\) supports and at most \(\lfloor t/r\rfloor\le n\) elements on the \(t\) deleted coordinates. This polynomial correction disappears in the exponent; radius zero is injective. Thus all closed branches in the stated envelope are retained.

\section{Improvement at a Conditional Support Boundary}
For a finite ensemble with mean \(P\), let \(\Pi_x\) project onto \(\supp\sigma_x\), and put
\(\kappa=q^{-1}\sum_x\Tr[(I-\Pi_x)P]\).

\begin{lemma}\label{lem:activation}
Every cyclic channel with \(0<p<\theta\) and \(\kappa>0\) has a cyclic replacement with strictly smaller error and information, using at most \(q\) additional dimensions.
\end{lemma}
\begin{proof}
Write \(h(a)=-a\ln a-(1-a)\ln(1-a)\) and
\(J_\varepsilon(a)=h(\varepsilon a)-\varepsilon h(a)\).
Encoding a binary label by \(\sigma_x\) and \(P\), with probabilities \(1-\varepsilon,\varepsilon\), and measuring
\(\{\Pi_x,I-\Pi_x\}\), gives information \(J_\varepsilon(\Tr[(I-\Pi_x)P])\). The Holevo measurement inequality~\cite{Holevo1973} therefore gives, for
\(\sigma_x^\varepsilon=(1-\varepsilon)\sigma_x+\varepsilon P\),
\begin{equation}\label{eq:entropyactivation}
\chi_\varepsilon\le(1-\varepsilon)\chi-J_\varepsilon(\kappa).
\end{equation}
Here Jensen applies because
\[
J_\varepsilon''(a)=
\frac{\varepsilon(1-\varepsilon)}{(1-a)(1-\varepsilon a)}>0,\qquad
J_\varepsilon(a)\ge\varepsilon a[\ln(1/\varepsilon)-1].
\]
The latter follows by expanding the two binary entropies and using
\((1-a)\ln(1-a)\ge-a\).

The mean stays \(P\). Expanding~\eqref{eq:cq} gives the exact error
\(p_\varepsilon=p+(2\varepsilon-\varepsilon^2)(\theta-p)\).
For \(0<\varepsilon<\min(1,\sqrt p)\), use the noiseless flag
\begin{align}
t_\varepsilon&=\frac{(2\varepsilon-\varepsilon^2)(\theta-p)+\varepsilon^2}
{p_\varepsilon},\notag\\
\tau_x&=(1-t_\varepsilon)\sigma_x^\varepsilon
\oplus t_\varepsilon|x\rangle\langle x|.\label{eq:repair}
\end{align}
Then \(0<t_\varepsilon<1\), \(p_\tau=p-\varepsilon^2\), and
\[
\chi-\chi_\tau\ge\varepsilon\chi+
J_\varepsilon(\kappa)-t_\varepsilon\ln q>0
\]
for sufficiently small \(\varepsilon\), since \(t_\varepsilon=O(\varepsilon)\) and \(\kappa>0\).
\end{proof}

For a classical generator~\eqref{eq:lift}, \(\kappa=0\) implies commuting outputs. If some column \(v_j=0\), a positive-trace mean block is absent from an output support, giving \(\kappa>0\). Otherwise zero defect requires
\(\operatorname{range}G\subseteq\operatorname{span}v_j\) for every \(j\), hence \(\rank G=1\). All blocks are then multiples of the same projection and commute.

\section{The Full-Rank Mixed-Channel Generators}\label{sec:MC}
The unmasked \(q\)-MC family of~\cite[Sec.~7]{AG2026}, in Fourier coordinates, is the phase orbit of
\begin{equation}\label{eq:MC}
\begin{gathered}
A_t=(1-t)P+tww^*,\qquad w_j=\sqrt{P_{jj}},\\
P=\diag(1-\alpha,\alpha/k,\ldots,\alpha/k),
\end{gathered}
\end{equation}
with \(0\le\alpha\le\theta\), \(0\le t\le1\), and phase unitary
\(U=\diag(1,\zeta,\ldots,\zeta^{q-1})\), \(\zeta=e^{2\pi i/q}\).
Let \(\mathfrak M_q\) include this entire compact domain.
For \(0<\alpha<\theta\), use
\begin{gather*}
r=\sqrt{k(1-\alpha)/\alpha}>1,\qquad Z=r^2+k,\\
L=2r+k-1,\qquad A=(r-1)(r+k).
\end{gather*}
The diagonal is \((b,B,\ldots,B)=(r^2/Z,1/Z,\ldots,1/Z)\), and the two coherence classes are \(c=tr/Z\) and \(d=t/Z\).
Set \(u=1-t\), \(v=1+kt\), and \(w_0=1+(k-1)t\). The block
\[
K=\begin{pmatrix}r^2&\sqrt{k}tr\\\sqrt{k}tr&w_0\end{pmatrix}
\]
has eigenvalues \(\nu_+>\nu_->0\), with \(\nu_- = r^2uv/\nu_+\).
Direct diagonalization gives
\begin{align}
p&=\theta(1-t^2L/Z),\label{eq:MCerror}\\
\chi_{\rm MC}&=Z^{-1}\big[\nu_+\ln\nu_++\nu_-\ln\nu_-
+(k-1)u\ln u-2r^2\ln r\big].\label{eq:MCchi}
\end{align}

\subsection{Two independent fixed-error directions}
Define
\begin{align}
m&=\frac{\ln(\nu_+/\nu_-)}{\nu_+-\nu_-},\notag\\
L_{00}&=\ln\nu_-+m(r^2-\nu_-),&
L_{++}&=\ln\nu_-+m(w_0-\nu_-),\label{eq:logs}\\
D_\perp&=\ln u-L_{++}+ktrm.\notag
\end{align}
Keep the diagonal fixed and vary the two positive coherences independently. The seed is positive definite exactly when
\[
B-d>0,\qquad b[B+(k-1)d]-kc^2>0,
\]
and its error is
\begin{equation}\label{eq:twocoherence}
p=\theta\left(1-\frac{2c^2}{\sqrt{bB}}-\frac{(k-1)d^2}{B}\right).
\end{equation}
At the original point use
\[
d_h=d+h,\qquad c_h=\sqrt{c^2-(k-1)r(d_h^2-d^2)/2}.
\]
This preserves error exactly and has \(c'_0=-(k-1)/2\). Matrix differentiation gives
\begin{equation}\label{eq:transverse}
\left.\frac{d\chi}{dh}\right|_{0}=-(k-1)D_\perp.
\end{equation}
Indeed the first coherence of \(\ln A_t\) is \(trm\), while the other is \((L_{++}-\ln u)/k\); contraction with the two perturbations proves~\eqref{eq:transverse}.
If \(D_\perp\ne0\), choose the decreasing sign of \(h\), then slightly multiply both new coherences by a factor greater than one. Positive definiteness is open, error decreases strictly by~\eqref{eq:twocoherence}, and continuity preserves the information gain.

Along a fixed-error curve inside the original family, \(t'=tA/(ZL)\). Set
\[
D_\parallel=\partial_r\chi_{\rm MC}+\frac{tA}{ZL}\partial_t\chi_{\rm MC}.
\]
Using~\eqref{eq:MCchi} gives
\begin{align}
\partial_r\chi_{\rm MC}
&=\frac{2r}{Z}(L_{00}-2\ln r+kt^2m-\chi_{\rm MC}),\notag\\
\partial_t\chi_{\rm MC}
&=\frac{2ktr^2m+(k-1)(L_{++}-\ln u)}Z.\label{eq:derivatives}
\end{align}
Introduce
\begin{align}
T_0&=\ln(v/u)-tm(qr+tA/2),\notag\\
C_0&=\frac{q+(k-1)t}{k}-\frac{(k-1)tA}{2krL}.\notag
\end{align}
Then
\begin{equation}\label{eq:coupled}
D_\parallel=\frac{2kr}{Z^2}(T_0+C_0D_\perp).
\end{equation}
To verify this identity substitute
\(L_{00}+L_{++}=\ln(r^2uv)\) and
\(L_{++}=\ln u+ktrm-D_\perp\) into~\eqref{eq:derivatives}.
If \(D_\parallel\ne0\), move in the decreasing direction along the fixed-error curve, then increase \(t\) slightly. This also strictly improves both coordinates. Only simultaneous zeros can remain.

\section{There Are Only Finitely Many Exceptional Points}
\begin{lemma}\label{lem:zerocurve}
For each \(r>1\), \(D_\perp(r,t)\) has exactly one zero \(t_*(r)\in(0,1)\), is negative below it and positive above it, and \(t_*\) is real analytic.
\end{lemma}
\begin{proof}
Put \(a=t/(1-t)\) and
\begin{gather*}
H_0(y)=(y+r^2)(y+1),\\
H_1(y)=(r^2+k)y+qr^2,\qquad H_a=H_0+aH_1.
\end{gather*}
The matrix-log resolvent, followed by \(z=(1-t)y\), gives
\begin{equation}\label{eq:resolvent}
D_\perp=ka(r-1)I(a,r),\qquad
I(a,r)=\int_0^\infty\frac{y-r}{(y+1)H_a(y)}\,dy.
\end{equation}
For example the unscaled integrand for \(D_\perp/[kt(r-1)]\) is
\((z-r(1-t))/[(z+1-t)\det(K+zI)]\), which verifies the normalization.

Normalize \(W_a(y)=1/[(y+1)H_a(y)]\) to a probability density. Its first moment is finite and \(I=0\) means \(\mathbb E_aY=r\).
For \(h_0=H_0/H_1\),
\[
H_1^2h'_0=(r^2+k)y^2+2qr^2y+r^2(kr^2+1)>0.
\]
Thus
\[
\frac{d}{da}\mathbb E_aY
=\operatorname{Cov}_a\left(Y,-\frac1{h_0(Y)+a}\right)>0.
\]
The two functions increase strictly; the bounded score and finite first moment justify differentiation and the strict covariance inequality.
The spectral formulas give
\begin{align}
\lim_{t\downarrow0}\frac{D_\perp}{t}
&=k\left(\frac{2r\ln r}{r^2-1}-1\right)<0,\notag\\
\lim_{t\uparrow1}\frac{D_\perp}{\ln(1/(1-t))}
&=\frac{k(r-1)}{r^2+k}>0.\label{eq:ends}
\end{align}
For the first sign use \(r-r^{-1}>2\ln r\); the second uses
\(\nu_+\to Z\), \(\nu_-/(1-t)\to qr^2/Z\).
Existence and uniqueness follow. At a zero \(I_a>0\), so the analytic implicit-function theorem applies.
\end{proof}

\begin{lemma}\label{lem:finite}
The set
\(\mathcal R_q=\{r>1:D_\parallel(r,t_*(r))=0\}\) is finite and nonempty.
\end{lemma}
\begin{proof}
First examine \(r=1+\eta\). Differentiation of the rational integral in~\eqref{eq:resolvent} gives
\[
I(0,1)=0,\qquad I_r(0,1)=-1/6,\qquad I_a(0,1)=q/6.
\]
Locally it extends analytically to negative \(a\): \(H_0/H_1\ge1/q\), so \(|a|<1/(2q)\) keeps \(H_a\ge H_0/2\).
Consequently \(t_*(1+\eta)=\eta/q+O_q(\eta^2)\).
For \(t=c\eta+d\eta^2\), expansion of \(\ln(I+E)\) in the two-dimensional block gives
\begin{align}
D_\perp&=\frac{kc(qc-1)}6\eta^3+O(\eta^4),\notag\\
T_0&=-\frac{cq(cq+c-1)}6\eta^3+O(\eta^4).\label{eq:cubic}
\end{align}
The cubic coefficients do not depend on \(d\). Substituting \(c=1/q\) into~\eqref{eq:coupled} yields
\begin{equation}\label{eq:nearone}
D_\parallel(1+\eta,t_*(1+\eta))
=-\frac{k}{3q^3}\eta^3+O_q(\eta^4)<0.
\end{equation}

At the other endpoint, uniformly for \(0<t<1\),
\(r^2\le\nu_+\le r^2+k\) and \(0<\nu_-\le k\).
For any fixed \(t_0\in(0,1)\), \(\nu_-/(1-t_0)\to1+kt_0\) and \(m=O(\ln r/r^2)\), so the stable expression
\[
D_\perp=\ln((1-t)/\nu_-)+m(ktr-w_0+\nu_-)
\]
tends to \(-\ln(1+kt_0)<0\).
Lemma~\ref{lem:zerocurve} forces \(t_*(r)\to1\). Along this curve \(\nu_-/(1-t_*)\to q\), and its root equation is
\[
m(kt_*r-w_0+\nu_-)=\ln(\nu_-/(1-t_*)).
\]
Therefore \(rm\to\ln q/k\). Since \((\nu_+-\nu_-)/r^2\to1\),
\[
\frac{\ln(1/(1-t_*))}{r}\longrightarrow\frac{\ln q}{k},
\qquad \frac{T_0}{r}\longrightarrow\frac{\ln q}{2k}.
\]
The first limit also follows by subtracting \(\ln\nu_+\) from \(\ln(\nu_+/\nu_-)\) and adding \(\ln(\nu_-/(1-t_*))\).
Equation~\eqref{eq:coupled} now gives
\[
r^2D_\parallel(r,t_*(r))\longrightarrow\ln q>0.
\]
The function is analytic and has opposite strict signs at its two ends. Its zeros lie in a compact interior interval; being isolated there, they are finite, and the sign change gives at least one.
\end{proof}

\section{The Joint Optimizer Must Be Improvable}
Call a generator improvable if some finite cyclic replacement has strictly smaller \(p\) and \(\chi\). Every generator is improvable unless it is commuting, is a zero-error or zero-information endpoint, or is one of the finite points of Lemma~\ref{lem:finite}.
For the classical family this follows from Lemma~\ref{lem:activation} and the rank-one argument following it. For \(q\)-MC, \(\alpha=0\) or \(t=0\) is uninformative; at \(\alpha=\theta\) the pure orbit is an orthonormal Fourier basis, so its mixtures commute. At \(t=1\), \(0<\alpha<\theta\), the output is pure and
\(\kappa=1-\Tr(PA_1)>0\), so Lemma~\ref{lem:activation} applies. All remaining points have \(r>1,0<t<1\), and the two coherence arguments apply unless both derivatives vanish.

We need only a simple slope bound for commuting channels. In a common eigenbasis let \(a_y\) be the output distribution and \(Q_y\) the input posterior under the uniform prior \(u_q\). Then
\[
\theta-p=\sum_y a_y\|Q_y-u_q\|_2^2,\qquad
\chi=\sum_y a_yD(Q_y\Vert u_q).
\]
Pinsker's inequality gives
\begin{equation}\label{eq:commutingslope}
\chi\ge\tfrac12\sum_y a_y\|Q_y-u_q\|_1^2
\ge\tfrac12(\theta-p).
\end{equation}
Every exceptional interior point has \(0<p<\theta\) and positive information. Hence the finite set defines
\begin{equation}\label{eq:beta}
\beta_q=\min\left\{\frac12,\
\min_{r\in\mathcal R_q}
\frac{\chi_{\rm MC}(r,t_*(r))}{\theta-p(r,t_*(r))}\right\}>0.
\end{equation}
A zero-error channel has \(\chi=\ln q\) by the measurement-information inequality, and a zero-information channel has \(p=\theta\). Consequently every nonimprovable generator satisfies
\begin{equation}\label{eq:slope}
\chi\ge\beta_q(\theta-p).
\end{equation}

\begin{proof}[Proof of Theorem~\ref{thm:main}]
The pair set in~\eqref{eq:benchmark} is compact and contains \((0,1)\) and \((\theta,0)\). Information zero forces every output to equal its mean, hence \(p=\theta\). Conversely, with \(Z_x=P^{-1/4}\sigma_xP^{-1/4}\),
\[
\theta-p=q^{-2}\sum_x\Tr[(Z_x-P^{1/2})^2],
\]
so \(p=\theta\) also forces equal outputs. Thus \(B_q\) is positive below \(\theta\), convex and nonincreasing. Its minimum has actual error \(\delta\), since otherwise mixing with \((\theta,0)\) would strictly lower its value. A supporting line shows that at most two original generators attain it, as in Lemma~\ref{lem:closure}. Finite convexity gives continuity in the interior.

Write \(\delta=\theta(1-s)\). The included first-LP source satisfies
\[
\gamma_q(\delta)=\frac{ks^2}{(\sqrt{1+ks}+\sqrt{1-s})^2}\le ks^2.
\]
Using \(h(x)+x\ln k\le x\ln(ek/x)\) yields, for
\(0<s\le(\beta_q/(4q))^2\),
\begin{equation}\label{eq:tail}
(\ln q)B_q(\delta)\le ks^2[1+2\ln(1/s)]
<\beta_q(\theta-\delta).
\end{equation}
Indeed division by \(\theta s\) gives \(qs[1+2\ln(1/s)]\), increasing in this range. At its upper endpoint, put \(x=\beta_q/(4q)\le1/24\); this ratio is
\((\beta_q/4)x[1+4\ln(1/x)]<\beta_q/4\), since the increasing factor is less than \(17/24\).

Let an attaining mixture have weights \(\lambda,1-\lambda\). If all positive-weight components were nonimprovable, their affine pairs and~\eqref{eq:slope} would contradict~\eqref{eq:tail}. Replace an improvable component of positive weight by its improving channel. The new mixture has strictly smaller error and information, proving~\eqref{eq:strict} via~\eqref{eq:criterion}.
Each original component has dimension at most \(q^2\); a boundary repair adds at most \(q\), and a full-rank coherence improvement uses dimension \(q\). The two-component mixture therefore has dimension at most \(2q^2+q\).

At an interior \(\delta_0\), retain this single channel. Its strict error margin persists nearby, and continuity of \(B_q\) preserves a positive information gap on a smaller compact neighborhood. All channel parameters are chosen before taking the blocklength limit.
\end{proof}

\begin{thebibliography}{20}


\bibitem{AG2026}
O. Alrabiah and V. Guruswami, ``Binary code rate bounds via classical--quantum channels,'' Aug. 2026. \arxiv{2608.09347v1}. doi: \doi{10.48550/arXiv.2608.09347}.

\bibitem{Aaltonen1990}
M. Aaltonen, ``A new upper bound on nonbinary block codes,'' \emph{Discrete Math.}, vol.~83, nos.~2--3, pp.~139--160, 1990. doi: \doi{10.1016/0012-365X(90)90002-Y}.

\bibitem{BHL2007}
Y. Ben-Haim and S. Litsyn, ``A new upper bound on the rate of non-binary codes,'' \emph{Adv. Math. Commun.}, vol.~1, no.~1, pp.~83--92, 2007. doi: \doi{10.3934/amc.2007.1.83}.

\bibitem{LL1998}
T. Laihonen and S. Litsyn, ``On upper bounds for minimum distance and covering radius of non-binary codes,'' \emph{Des. Codes Cryptogr.}, vol.~14, pp.~71--80, 1998. doi: \doi{10.1023/A:1008260505585}.

\bibitem{Holevo1973}
A. S. Holevo, ``Bounds for the quantity of information transmitted by a quantum communication channel,'' \emph{Probl. Inf. Transm.}, vol.~9, no.~3, pp.~177--183, 1973. \href{https://www.mathnet.ru/eng/ppi903}{Journal archive}.
\end{thebibliography}

\end{document}
